\documentclass[../DoAn.tex]{subfiles}
\begin{document}

Chương 2 đã xác định phạm vi chức năng và các yêu cầu chính của trò chơi giải đố từ vựng dạng ô chữ. Chương 3 đã trình bày nền tảng công nghệ và các vấn đề lý thuyết cần thiết như Godot, GDScript, xử lý chuỗi tiếng Việt, sinh bảng ô chữ và lưu dữ liệu cục bộ. Trên cơ sở đó, chương này trình bày cách chuyển các yêu cầu thành thiết kế cụ thể, mô tả kiến trúc, luồng hoạt động, module, dữ liệu lưu trữ, giao diện, kết quả xây dựng nguyên mẫu và khung kiểm thử.

Nội dung chương tập trung vào hệ thống đã triển khai trong project Godot hiện tại. Các phần liên quan đến kiểm thử được trình bày theo hướng khung kiểm thử và tiêu chí ghi nhận, không tự kết luận Passed/Failed khi chưa có dữ liệu chạy kiểm thử chính thức.

\section{Thiết kế kiến trúc hệ thống}
\label{section:4.1}

Ứng dụng được thiết kế theo hướng phân tách trách nhiệm giữa điều phối ứng dụng, giao diện, trạng thái trò chơi, sinh nội dung, xử lý tiếng Việt và lưu trữ cục bộ. Do Godot tổ chức chương trình theo scene và node, thiết kế không áp dụng mô hình MVC theo nghĩa chặt chẽ. Tuy nhiên, hệ thống vẫn tách rõ phần hiển thị, phần trạng thái, phần sinh bảng và phần lưu dữ liệu để giảm phụ thuộc trực tiếp giữa các thành phần.

Thành phần điều phối ứng dụng nằm chủ yếu trong \code{main.gd}. Script này quản lý luồng từ menu chính, chọn tài khoản, chọn chủ đề, chọn màn và chuyển sang màn chơi. Giao diện màn chơi do \code{UIController} điều khiển, trong khi trạng thái động của một lượt chơi được quản lý bởi \code{GameState}. Dữ liệu chủ đề và danh sách từ được chuẩn bị bởi \code{ThemedLevelProvider}; quá trình đặt từ lên lưới do \code{BoardGenerator} thực hiện. Các vấn đề riêng của tiếng Việt được gom vào \code{WordEntry}, \code{VietnameseNormalizer} và \code{DiacriticHelper}. Dữ liệu tiến độ được đọc ghi thông qua \code{SaveManager} với tệp \code{user://records.json}.

\begin{figure}[H]
\centering
\resizebox{0.95\textwidth}{!}{\begin{tikzpicture}[
    font=\small,
    block/.style={draw, rounded corners=5pt, align=center, minimum width=3.2cm, minimum height=1.0cm, fill=gray!8},
    store/.style={draw, rounded corners=5pt, align=center, minimum width=3.2cm, minimum height=1.0cm, fill=yellow!12},
    actor/.style={draw, rounded corners=5pt, align=center, minimum width=2.8cm, minimum height=0.9cm, fill=blue!8},
    arrow/.style={-{Stealth[length=2mm]}, thick}
]
\node[actor] (player) at (0,0) {Người chơi};
\node[block] (ui) at (4.0,0) {Giao diện\\Godot};
\node[block] (main) at (8.0,0) {Điều phối\\ứng dụng};
\node[block] (state) at (8.0,-2.0) {Trạng thái\\trò chơi};
\node[block] (board) at (4.0,-4.0) {Sinh bảng\\ô chữ};
\node[block] (vietnamese) at (8.0,-4.0) {Xử lý\\tiếng Việt};
\node[store] (topic) at (0,-4.0) {Dữ liệu\\chủ đề};
\node[store] (save) at (12.0,-2.0) {Lưu tiến độ\\JSON};

\draw[arrow] (player) -- (ui);
\draw[arrow] (ui) -- (main);
\draw[arrow] (main) -- (state);
\draw[arrow] (state) -- (vietnamese);
\draw[arrow] (state) -- (board);
\draw[arrow] (topic) -- (board);
\draw[arrow] (main) -- (save);
\draw[arrow] (state) -- (save);
\draw[arrow] (state) -- (ui);
\end{tikzpicture}
}
\caption{Kiến trúc tổng quan của hệ thống}
\label{fig:kien-truc-he-thong}
\end{figure}

\begin{table}[H]
\centering
\small
\begin{tabularx}{\textwidth}{|>{\RaggedRight\arraybackslash}p{3.0cm}|>{\RaggedRight\arraybackslash}p{4.2cm}|>{\RaggedRight\arraybackslash}X|}
\hline
\textbf{Nhóm thành phần} & \textbf{Script/scene chính} & \textbf{Trách nhiệm} \\ \hline
Điều phối ứng dụng & \code{main.gd}, \code{Main.tscn} & Quản lý menu, tài khoản, chủ đề, màn chơi và chuyển đổi giữa menu với màn chơi. \\ \hline
Giao diện & \code{UIController}, \code{StartMenu.tscn}, \code{GameBoard.tscn}, \code{Tile.tscn} & Hiển thị bảng ô chữ, danh sách gợi ý, thông tin thời gian, sao, điểm, hộp thoại tạm dừng và hộp thoại hoàn thành. \\ \hline
Trạng thái trò chơi & \code{GameState}, \code{InputController}, \code{TileButton} & Lưu trạng thái một lượt chơi, từ đang chọn, lịch sử đoán, điểm, nhiệm vụ, trạng thái các ô và trạng thái hoàn thành từ. \\ \hline
Sinh nội dung & \code{ThemedLevelProvider}, \code{BoardGenerator}, \code{ManualClues} & Tổ chức dữ liệu chủ đề, tạo danh sách từ, sinh gợi ý và đặt từ lên bảng ô chữ. \\ \hline
Xử lý tiếng Việt & \code{WordEntry}, \code{VietnameseNormalizer}, \code{DiacriticHelper} & Chuẩn hóa đáp án, bỏ khoảng trắng khi xếp bảng, tách ký tự tiếng Việt và lưu thông tin dấu. \\ \hline
Lưu trữ cục bộ & \code{SaveManager} & Đọc, ghi và chuẩn hóa dữ liệu tài khoản, sao, điểm cao nhất và tiến độ màn chơi. \\ \hline
\end{tabularx}
\caption{Phân nhóm thành phần trong kiến trúc ứng dụng}
\label{table:architecture_groups}
\end{table}

Cách phân tách trên giúp các thành phần có trách nhiệm rõ hơn. Ví dụ, \code{BoardGenerator} không cần biết bảng được hiển thị bằng node nào; nó chỉ trả về lưới và danh sách từ đã đặt. Tương tự, \code{GameState} không phụ thuộc trực tiếp vào scene giao diện mà chỉ cung cấp các thao tác như kiểm tra đáp án, mở từ, tính điểm và đánh giá nhiệm vụ. Giao diện dùng các kết quả đó để cập nhật hiển thị.

\section{Thiết kế luồng hoạt động}
\label{section:4.2}

Luồng hoạt động chính bắt đầu từ menu. Người chơi chọn hoặc tạo tài khoản cục bộ, bấm Chơi, chọn chủ đề, chọn màn, sau đó hệ thống khởi tạo trạng thái chơi và hiển thị màn chơi. Khi hoàn thành màn, hộp thoại kết quả xuất hiện và người chơi có thể chơi lại hoặc quay về menu/chọn màn. Hình~\ref{fig:luong-chuyen-man} mô tả luồng chuyển màn ở mức giao diện.

\begin{figure}[H]
\centering
\resizebox{0.92\textwidth}{!}{\begin{tikzpicture}[
    font=\small,
    screenstep/.style={draw, rounded corners=5pt, align=center, minimum width=3.3cm, minimum height=0.9cm, fill=gray!8},
    arrow/.style={-{Stealth[length=2mm]}, thick}
]
\node[screenstep] (menu) at (0,0) {StartMenu};
\node[screenstep] (profile) at (3.8,0) {Chọn/Tạo\\tài khoản};
\node[screenstep] (topic) at (7.6,0) {Chọn\\chủ đề};
\node[screenstep] (level) at (11.4,0) {Chọn\\màn};
\node[screenstep] (state) at (11.4,-2.0) {Khởi tạo\\GameState};
\node[screenstep] (board) at (7.6,-2.0) {GameBoard};
\node[screenstep] (result) at (3.8,-2.0) {Hộp thoại\\kết quả};
\node[screenstep] (back) at (0,-2.0) {Quay lại menu\\hoặc chọn màn};

\draw[arrow] (menu) -- (profile);
\draw[arrow] (profile) -- (topic);
\draw[arrow] (topic) -- (level);
\draw[arrow] (level) -- (state);
\draw[arrow] (state) -- (board);
\draw[arrow] (board) -- node[above]{hoàn thành} (result);
\draw[arrow] (result) -- (back);
\draw[arrow] (back) -- (menu);
\draw[arrow] (back.north east) .. controls (3.0,-0.9) and (8.0,-0.9) .. node[below]{chơi tiếp} (level.south west);
\end{tikzpicture}
}
\caption{Luồng chuyển màn trong ứng dụng}
\label{fig:luong-chuyen-man}
\end{figure}

Hai quy trình xử lý quan trọng nhất trong luồng này là khởi tạo một màn chơi và kiểm tra đáp án trong khi chơi. Các quy trình này được mô tả sau phần thiết kế module để người đọc có thể đối chiếu trình tự gọi hàm với vai trò của từng script.

\section{Thiết kế module và lớp}
\label{section:4.3}

Các script chính được thiết kế quanh dữ liệu và luồng chơi của một màn ô chữ. \code{WordEntry} giữ thông tin của một đáp án gồm từ gốc, dạng dùng để xếp bảng, gợi ý, danh sách ký tự và thông tin dấu. \code{BoardGenerator} nhận danh sách \code{WordEntry} và đặt các từ lên lưới. \code{GameState} quản lý trạng thái động sau khi bảng được sinh. \code{UIController} biến trạng thái thành giao diện và nhận thao tác từ người chơi. \code{SaveManager} xử lý dữ liệu lưu cục bộ.

\begin{figure}[H]
\centering
\resizebox{0.90\textwidth}{!}{\begin{tikzpicture}[
    font=\scriptsize,
    class/.style={draw, rounded corners=4pt, align=center, minimum width=2.8cm, minimum height=0.72cm, fill=gray!8},
    data/.style={draw, rounded corners=4pt, align=center, minimum width=2.8cm, minimum height=0.72cm, fill=blue!8},
    util/.style={draw, rounded corners=4pt, align=center, minimum width=2.8cm, minimum height=0.72cm, fill=yellow!12},
    arrow/.style={-{Stealth[length=2mm]}, thick}
]
\node[class] (main) at (5.8,3.0) {\code{main.gd}};
\node[class] (ui) at (2.8,1.5) {\code{UIController}};
\node[class] (provider) at (8.8,1.5) {\code{ThemedLevelProvider}};
\node[util] (save) at (11.8,3.0) {\code{SaveManager}};
\node[class] (input) at (0,0) {\code{InputController}};
\node[class] (tile) at (2.8,0) {\code{TileButton}};
\node[data] (state) at (5.8,0) {\code{GameState}};
\node[class] (board) at (8.8,0) {\code{BoardGenerator}};
\node[data] (word) at (5.8,-1.6) {\code{WordEntry}};
\node[util] (normalizer) at (2.8,-3.0) {\code{VietnameseNormalizer}};
\node[util] (diacritic) at (8.8,-3.0) {\code{DiacriticHelper}};

\draw[arrow] (main) -- (ui);
\draw[arrow] (main) -- (provider);
\draw[arrow] (main) -- (save);
\draw[arrow] (ui) -- (input);
\draw[arrow] (ui) -- (tile);
\draw[arrow] (ui) -- (state);
\draw[arrow] (provider) -- (board);
\draw[arrow] (provider) -- (word);
\draw[arrow] (state) -- (word);
\draw[arrow] (board) -- (word);
\draw[arrow] (word) -- (normalizer);
\draw[arrow] (word) -- (diacritic);
\draw[arrow] (normalizer) -- (diacritic);
\draw[arrow] (save) -- node[right]{JSON} ++(0,-1.6);
\end{tikzpicture}
}
\caption{Class diagram rút gọn của các thành phần chính}
\label{fig:class-diagram-rut-gon}
\end{figure}

\begin{table}[H]
\centering
\small
\begin{tabularx}{\textwidth}{|>{\RaggedRight\arraybackslash}p{3.6cm}|>{\RaggedRight\arraybackslash}p{4.8cm}|>{\RaggedRight\arraybackslash}X|}
\hline
\textbf{Script/lớp} & \textbf{Dữ liệu hoặc phụ thuộc chính} & \textbf{Chức năng tiêu biểu} \\ \hline
\code{main.gd} & \code{BoardGenerator}, \code{ThemedLevelProvider}, \code{SaveManager}, \code{UIController} & Điều phối menu, tài khoản, chủ đề, danh sách màn và khởi tạo \code{GameState}. \\ \hline
\code{UIController} & \code{GameState}, \code{InputController}, \code{TileButton}, các node trong \code{GameBoard.tscn} & Dựng bảng, tạo nút gợi ý, nhận đáp án, dùng gợi ý, cập nhật điểm/sao/thời gian và hiển thị kết quả. \\ \hline
\code{InputController} & \code{UIController}, \code{GameState}, vị trí ô và từ đang chọn & Lưu trạng thái chọn ô/chọn từ và hỗ trợ chuyển lựa chọn giữa các từ giao nhau. \\ \hline
\code{TileButton} & Dữ liệu ô trong bảng & Hiển thị một ô chữ, phân biệt ô rỗng, ô đang chọn, ô thuộc từ đang chọn và ô đã điền đúng. \\ \hline
\code{GameState} & Lưới, danh sách từ đã đặt, điểm, thời gian, nhiệm vụ, lịch sử đoán & Kiểm tra từ, mở từ, tính điểm, đếm tiến độ, quản lý nhiệm vụ và ghi nhận lịch sử đoán. \\ \hline
\code{BoardGenerator} & Danh sách \code{WordEntry}, kích thước bảng & Tạo bảng rỗng, tìm vị trí giao nhau, kiểm tra đặt từ hợp lệ, chấm điểm vị trí đặt và loại bỏ từ không liên kết. \\ \hline
\code{ThemedLevelProvider} & Danh sách chủ đề, \code{ManualClues}, \code{BoardGenerator} & Tạo chủ đề, chia màn, tạo \code{WordEntry}, sinh gợi ý và yêu cầu sinh bảng cho từng màn. \\ \hline
\code{WordEntry} & Từ gốc, dạng xếp bảng, gợi ý, danh sách ký tự & Chuẩn bị dữ liệu của một đáp án để vừa xếp bảng vừa kiểm tra đáp án. \\ \hline
\code{VietnameseNormalizer} & Chuỗi tiếng Việt & Chuẩn hóa chữ thường, bỏ khoảng trắng khi cần, tách ký tự và so khớp đáp án. \\ \hline
\code{DiacriticHelper} & Bảng ánh xạ dấu và nguyên âm tiếng Việt & Tách ký tự thành chữ cơ sở, biến thể nguyên âm, thanh điệu và hỗ trợ bỏ dấu. \\ \hline
\code{SaveManager} & \code{user://records.json} & Tạo/chọn/xóa tài khoản, lưu sao, điểm cao nhất, số lượt chơi và tiến độ từng màn. \\ \hline
\end{tabularx}
\caption{Mô tả các lớp và script chính}
\label{table:main_classes}
\end{table}

\begin{table}[H]
\centering
\small
\begin{tabularx}{\textwidth}{|>{\RaggedRight\arraybackslash}p{3.7cm}|>{\RaggedRight\arraybackslash}X|}
\hline
\textbf{Lớp/script} & \textbf{Thiết kế chi tiết} \\ \hline
\code{GameState} & Lớp trung tâm của trạng thái lượt chơi. Dữ liệu gồm lưới, từ đã đặt, số lần đoán, số đáp án đúng, điểm, thời gian, nhiệm vụ, lịch sử đoán, số gợi ý đã dùng và từ đang chọn. Lớp này không trực tiếp dựng giao diện. \\ \hline
\code{BoardGenerator} & Lớp thuật toán sinh bảng. Từ đầu tiên được đặt gần trung tâm; các từ sau ưu tiên vị trí có giao điểm, diện tích bao nhỏ và khoảng cách hợp lý tới trung tâm. Sau khi đặt, các từ không có giao điểm với phần còn lại được loại bỏ để tránh bảng rời rạc. \\ \hline
\code{ThemedLevelProvider} & Lớp chuẩn bị nội dung. Script này chứa danh sách chủ đề, tạo \code{WordEntry}, chia chủ đề thành các màn nhỏ và gọi \code{BoardGenerator} nhiều lần để tìm bảng đủ điều kiện chơi. \\ \hline
\code{UIController} & Lớp điều khiển màn chơi. Script này tạo các \code{TileButton}, tạo danh sách nút gợi ý, xử lý nhập đáp án, xử lý gợi ý bằng sao, cập nhật thời gian, điểm, sao và mở hộp thoại hoàn thành. \\ \hline
\code{SaveManager} & Lớp tiện ích tĩnh cho lưu trữ. Dữ liệu được chuẩn hóa khi đọc để tránh lỗi khi tệp thiếu trường hoặc còn cấu trúc cũ. \\ \hline
\end{tabularx}
\caption{Thiết kế chi tiết các lớp chủ đạo}
\label{table:class_design}
\end{table}

\subsection{Trình tự xử lý chính}

Khi một màn được chọn, \code{main.gd} yêu cầu \code{ThemedLevelProvider} tạo dữ liệu màn. Dữ liệu này bao gồm thông tin màn, danh sách \code{WordEntry} và bảng ô chữ do \code{BoardGenerator} sinh ra. Sau khi dữ liệu hợp lệ, \code{main.gd} tạo \code{GameState}, gán bảng và thông tin màn, rồi truyền trạng thái này cho \code{UIController} để dựng giao diện. Hình~\ref{fig:sequence-bat-dau-man} trình bày trình tự này.

\begin{figure}[H]
\centering
\resizebox{0.95\textwidth}{!}{\begin{tikzpicture}[
    font=\footnotesize,
    lifeline/.style={draw, rounded corners=3pt, align=center, minimum width=2.1cm, minimum height=0.65cm, fill=gray!8},
    msg/.style={-{Stealth[length=2mm]}, thick},
    returnmsg/.style={-{Stealth[length=2mm]}, thick, dashed}
]
\node[lifeline] (player) at (0,0) {Người chơi};
\node[lifeline] (main) at (2.8,0) {\code{main.gd}};
\node[lifeline] (provider) at (5.8,0) {ThemedLevel\\Provider};
\node[lifeline] (board) at (9.0,0) {Board\\Generator};
\node[lifeline] (state) at (12.0,0) {GameState};
\node[lifeline] (ui) at (15.0,0) {UI\\Controller};

\foreach \x in {0,2.8,5.8,9.0,12.0,15.0}
    \draw[dashed, gray] (\x,-0.45) -- (\x,-6.3);

\draw[msg] (0,-1.0) -- node[above]{chọn màn} (2.8,-1.0);
\draw[msg] (2.8,-1.7) -- node[above]{yêu cầu dữ liệu màn} (5.8,-1.7);
\draw[msg] (5.8,-2.4) -- node[above]{tạo \code{WordEntry}} (5.8,-2.9);
\draw[msg] (5.8,-3.4) -- node[above]{sinh bảng} (9.0,-3.4);
\draw[returnmsg] (9.0,-4.1) -- node[above]{bảng và từ đã đặt} (5.8,-4.1);
\draw[returnmsg] (5.8,-4.8) -- node[above]{level + board} (2.8,-4.8);
\draw[msg] (2.8,-5.5) -- node[above]{khởi tạo trạng thái} (12.0,-5.5);
\draw[msg] (2.8,-6.1) -- node[above]{dựng màn chơi} (15.0,-6.1);
\end{tikzpicture}
}
\caption{Sequence diagram quy trình bắt đầu màn chơi}
\label{fig:sequence-bat-dau-man}
\end{figure}

Trong quá trình chơi, người chơi chọn gợi ý và nhập đáp án vào ô nhập. \code{InputController} giữ thông tin từ đang chọn, còn \code{UIController} nhận nội dung nhập và gọi \code{VietnameseNormalizer} để so khớp với đáp án gốc. Nếu đúng, \code{GameState} mở từ, cộng điểm, cập nhật lịch sử đoán và kiểm tra trạng thái hoàn thành. Khi toàn bộ từ được hoàn thành, \code{SaveManager} ghi nhận điểm, sao, thời gian tốt nhất và trạng thái màn. Hình~\ref{fig:sequence-nhap-dap-an} mô tả luồng nhập đáp án và cập nhật trạng thái.

\begin{figure}[H]
\centering
\resizebox{0.95\textwidth}{!}{\begin{tikzpicture}[
    font=\footnotesize,
    lifeline/.style={draw, rounded corners=3pt, align=center, minimum width=2.2cm, minimum height=0.65cm, fill=gray!8},
    msg/.style={-{Stealth[length=2mm]}, thick},
    returnmsg/.style={-{Stealth[length=2mm]}, thick, dashed},
    note/.style={draw, rounded corners=3pt, align=left, fill=yellow!12, text width=3.1cm}
]
\node[lifeline] (player) at (0,0) {Người chơi};
\node[lifeline] (ui) at (3.0,0) {UI\\Controller};
\node[lifeline] (input) at (6.0,0) {Input\\Controller};
\node[lifeline] (state) at (9.0,0) {\code{GameState}};
\node[lifeline] (normalizer) at (12.0,0) {Vietnamese\\Normalizer};
\node[lifeline] (save) at (15.0,0) {Save\\Manager};

\foreach \x in {0,3.0,6.0,9.0,12.0,15.0}
    \draw[dashed, gray] (\x,-0.45) -- (\x,-7.5);

\draw[msg] (0,-1.0) -- node[above]{chọn gợi ý} (3.0,-1.0);
\draw[msg] (3.0,-1.7) -- node[above]{cập nhật từ đang chọn} (6.0,-1.7);
\draw[msg] (0,-2.4) -- node[above]{nhập đáp án} (3.0,-2.4);
\draw[msg] (3.0,-3.1) -- node[above]{gửi đáp án} (9.0,-3.1);
\draw[msg] (9.0,-3.8) -- node[above]{chuẩn hóa/so khớp} (12.0,-3.8);
\draw[returnmsg] (12.0,-4.5) -- node[above]{đúng/sai} (9.0,-4.5);
\node[note] at (10.8,-5.35) {Nếu đúng: mở từ, cộng điểm, cập nhật lịch sử và trạng thái.};
\draw[msg] (9.0,-6.55) -- (15.0,-6.55);
\draw[returnmsg] (15.0,-7.15) -- node[above]{đã ghi tiến độ} (9.0,-7.15);
\end{tikzpicture}
}
\caption{Sequence diagram quy trình nhập đáp án và cập nhật trạng thái}
\label{fig:sequence-nhap-dap-an}
\end{figure}

\section{Thiết kế dữ liệu lưu trữ}
\label{section:4.4}

Ứng dụng sử dụng lưu trữ cục bộ thay vì cơ sở dữ liệu hoặc máy chủ. Dữ liệu được ghi vào \code{user://records.json}, là vùng lưu dữ liệu riêng của Godot trên máy người chơi. Cách lưu này phù hợp với phạm vi trò chơi một người chơi, không yêu cầu đăng nhập trực tuyến và không cần đồng bộ qua mạng.

\begin{figure}[H]
\centering
\resizebox{0.80\textwidth}{!}{\begin{tikzpicture}[
    font=\small,
    node/.style={draw, rounded corners=4pt, align=left, minimum width=3.0cm, minimum height=0.65cm, fill=gray!8},
    key/.style={draw, rounded corners=4pt, align=left, minimum width=3.0cm, minimum height=0.65cm, fill=blue!8},
    arrow/.style={-{Stealth[length=2mm]}, thick}
]
\node[key] (root) at (0,0) {\code{records.json}};
\node[node] (current) at (4.0,0.8) {\code{current_profile_id}};
\node[node] (profiles) at (4.0,-0.8) {\code{profiles}};
\node[node] (profile) at (8.0,-0.8) {\code{profile_id}};
\node[node] (name) at (12.0,1.6) {\code{name}};
\node[node] (records) at (12.0,0.4) {\code{records}};
\node[node] (score) at (16.0,2.7) {\code{best_score}};
\node[node] (stars) at (16.0,1.7) {\code{total_stars}};
\node[node] (earned) at (16.0,0.7) {\code{earned_stars}};
\node[node] (spent) at (16.0,-0.3) {\code{spent_stars}};
\node[node] (played) at (16.0,-1.3) {\code{games_played}};
\node[node] (levels) at (16.0,-2.3) {\code{levels}};
\node[node] (levelid) at (20.0,-2.3) {\code{level_id}};
\node[node] (completed) at (24.0,-0.8) {\code{completed}};
\node[node] (levelstars) at (24.0,-1.8) {\code{stars}};
\node[node] (bestscore) at (24.0,-2.8) {\code{best_score}};
\node[node] (besttime) at (24.0,-3.8) {\code{best_time}};

\draw[arrow] (root) -- (current);
\draw[arrow] (root) -- (profiles);
\draw[arrow] (profiles) -- (profile);
\draw[arrow] (profile) -- (name);
\draw[arrow] (profile) -- (records);
\draw[arrow] (records) -- (score);
\draw[arrow] (records) -- (stars);
\draw[arrow] (records) -- (earned);
\draw[arrow] (records) -- (spent);
\draw[arrow] (records) -- (played);
\draw[arrow] (records) -- (levels);
\draw[arrow] (levels) -- (levelid);
\draw[arrow] (levelid) -- (completed);
\draw[arrow] (levelid) -- (levelstars);
\draw[arrow] (levelid) -- (bestscore);
\draw[arrow] (levelid) -- (besttime);
\end{tikzpicture}
}
\caption{Cấu trúc dữ liệu lưu tiến độ cục bộ}
\label{fig:cau-truc-json-luu-tru}
\end{figure}

\begin{table}[H]
\centering
\small
\begin{tabularx}{\textwidth}{|>{\RaggedRight\arraybackslash}p{3.2cm}|>{\RaggedRight\arraybackslash}p{3.5cm}|>{\RaggedRight\arraybackslash}X|}
\hline
\textbf{Trường} & \textbf{Vị trí} & \textbf{Ý nghĩa} \\ \hline
\code{current_profile_id} & Gốc tệp & Mã tài khoản đang được chọn khi mở ứng dụng. \\ \hline
\code{profiles} & Gốc tệp & Tập các hồ sơ người chơi cục bộ, mỗi hồ sơ có một mã riêng. \\ \hline
\code{name} & Trong từng \code{profile_id} & Tên hiển thị của người chơi. \\ \hline
\code{best_score} & Trong \code{records} của hồ sơ & Điểm cao nhất từng ghi nhận ở tài khoản hiện tại. \\ \hline
\code{total_stars} & Trong \code{records} của hồ sơ & Số sao hiện còn, dùng để đổi gợi ý. \\ \hline
\code{earned_stars} & Trong \code{records} của hồ sơ & Tổng số sao đã nhận từ nhiệm vụ hoàn thành màn. \\ \hline
\code{spent_stars} & Trong \code{records} của hồ sơ & Tổng số sao đã dùng cho gợi ý. \\ \hline
\code{games_played} & Trong \code{records} của hồ sơ & Số lượt màn đã được ghi nhận sau khi hoàn thành. \\ \hline
\code{levels} & Trong \code{records} của hồ sơ & Tiến độ theo từng màn, gồm trạng thái hoàn thành, sao, điểm tốt nhất và thời gian tốt nhất. \\ \hline
\end{tabularx}
\caption{Các trường dữ liệu chính trong tệp lưu tiến độ}
\label{table:save_fields}
\end{table}

\code{SaveManager} có nhiệm vụ chuẩn hóa dữ liệu khi đọc. Nếu tệp lưu chưa tồn tại, hệ thống tạo cấu trúc mặc định với tài khoản \code{player}. Nếu dữ liệu cũ không có cấu trúc \code{profiles}, script chuyển dữ liệu về dạng hồ sơ mặc định để giữ khả năng đọc lại tiến độ cũ. Cách này giúp ứng dụng giảm khả năng dừng đột ngột khi tệp lưu thiếu trường.

\section{Thiết kế giao diện}
\label{section:4.5}

Giao diện gồm hai scene chính: \code{StartMenu.tscn} cho menu và \code{GameBoard.tscn} cho màn chơi. Trong menu, người chơi có thể chọn hoặc quản lý tài khoản, vào luồng chọn chủ đề và mở danh sách màn. Một số popup như quản lý tài khoản, thêm tài khoản, xác nhận xóa tài khoản, hướng dẫn và chọn màn được tạo bằng mã trong \code{main.gd} để nội dung có thể cập nhật theo trạng thái hiện tại.

Màn chơi được chia thành vùng bảng ô chữ và vùng điều khiển. Bảng dùng \code{GridContainer} để chứa các \code{TileButton}; vùng điều khiển chứa danh sách số gợi ý, ô nhập đáp án, nút đoán, nút gợi ý, trạng thái điểm/sao/thời gian và hộp thoại hoàn thành. Script \code{UIController} có hàm tính kích thước ô theo kích thước viewport, giới hạn bởi \code{MIN_TILE_SIZE} và \code{MAX_TILE_SIZE}, nhằm giảm nguy cơ bảng tràn khỏi vùng chơi.

\begin{table}[H]
\centering
\small
\begin{tabularx}{\textwidth}{|>{\RaggedRight\arraybackslash}p{3.2cm}|>{\RaggedRight\arraybackslash}p{4.4cm}|>{\RaggedRight\arraybackslash}X|}
\hline
\textbf{Màn hình} & \textbf{Thành phần chính} & \textbf{Mục đích thiết kế} \\ \hline
Menu chính & Tiêu đề, chọn tài khoản, nút Chơi, nút Hướng dẫn, nút Thoát & Cho người chơi bắt đầu nhanh và biết tài khoản đang dùng. \\ \hline
Chọn tài khoản & Popup quản lý tài khoản, thêm tài khoản, xóa tài khoản & Tách tiến độ của nhiều người chơi trên cùng máy. \\ \hline
Chọn chủ đề & Danh sách chủ đề kèm tiến độ hoàn thành & Giúp người chơi chọn nhóm từ vựng theo nội dung. \\ \hline
Chọn màn & Popup danh sách màn, số từ, số sao và trạng thái hoàn thành & Cho phép chọn màn cụ thể và xem nhanh tiến độ từng màn. \\ \hline
Màn chơi & Bảng ô chữ, gợi ý, ô nhập, nút đoán, nút gợi ý, thời gian và sao & Tập trung vào thao tác giải ô chữ và phản hồi đúng/sai. \\ \hline
Kết quả & Hộp thoại hoàn thành, điểm, sao, tỉ lệ đúng, thời gian, nhiệm vụ & Tổng kết lượt chơi và cho phép chơi lại hoặc quay về menu. \\ \hline
\end{tabularx}
\caption{Thiết kế các màn hình chính}
\label{table:ui_screens}
\end{table}

% TODO: Chèn ảnh màn hình menu chính sau khi có file ảnh thật.
% TODO: Chèn ảnh màn hình chọn chủ đề/chọn màn sau khi có file ảnh thật.
% TODO: Chèn ảnh màn hình chơi chính sau khi có file ảnh thật.
% TODO: Chèn ảnh hộp thoại kết quả sau khi có file ảnh thật.

Trong project hiện tại chưa có ảnh chụp màn hình game riêng để đưa vào báo cáo. Vì vậy, báo cáo chưa tạo các khung ảnh trống để tránh làm PDF xuất hiện placeholder không có nội dung.

\section{Xây dựng ứng dụng}
\label{section:4.6}

Ứng dụng được xây dựng bằng Godot 4.x và GDScript. Cấu trúc mã nguồn trong project hiện tại gồm thư mục \code{scripts} cho logic và thư mục \code{scenes} cho scene Godot. Thống kê được lấy trực tiếp từ các file trong project tại thời điểm chỉnh sửa báo cáo.

\begin{table}[H]
\centering
\small
\begin{tabularx}{\textwidth}{|>{\RaggedRight\arraybackslash}p{3.4cm}|>{\RaggedRight\arraybackslash}p{4.0cm}|>{\RaggedRight\arraybackslash}X|}
\hline
\textbf{Nhóm} & \textbf{Thành phần} & \textbf{Ghi chú sử dụng} \\ \hline
Game engine & Godot 4.x & Dùng để tổ chức scene, node, signal, giao diện 2D và chạy project cục bộ. \\ \hline
Ngôn ngữ & GDScript & Dùng cho toàn bộ logic điều phối, sinh bảng, trạng thái, giao diện và lưu dữ liệu. \\ \hline
Mã nguồn logic & 13 file \code{.gd}, 4538 dòng & Thống kê từ thư mục \code{scripts}. Không tính file \code{.uid}. \\ \hline
Scene & 4 file \code{.tscn}, 368 dòng & Gồm \code{Main.tscn}, \code{StartMenu.tscn}, \code{GameBoard.tscn}, \code{Tile.tscn}. \\ \hline
Lưu dữ liệu & \code{user://records.json} & Lưu tài khoản, sao, điểm, số lượt chơi và tiến độ màn. \\ \hline
\end{tabularx}
\caption{Danh sách công cụ và cấu trúc mã nguồn sử dụng}
\label{table:tools_source}
\end{table}

\begin{table}[H]
\centering
\small
\begin{tabularx}{\textwidth}{|>{\RaggedRight\arraybackslash}p{4.0cm}|>{\RaggedRight\arraybackslash}p{2.0cm}|>{\RaggedRight\arraybackslash}X|}
\hline
\textbf{Nhóm file} & \textbf{Số lượng} & \textbf{Vai trò} \\ \hline
\code{scripts/*.gd} & 13 file & Chứa logic trò chơi, giao diện, sinh bảng, xử lý tiếng Việt, chủ đề và lưu tiến độ. \\ \hline
\code{scenes/*.tscn} & 4 file & Chứa cấu trúc node của menu, màn chơi, scene chính và ô chữ. \\ \hline
\code{project.godot} & 1 file & Cấu hình project Godot và scene khởi động. \\ \hline
\end{tabularx}
\caption{Thống kê cấu trúc source hiện có}
\label{table:source_statistics}
\end{table}

Kết quả xây dựng nguyên mẫu gồm luồng chọn tài khoản, chọn chủ đề, chọn màn, sinh bảng ô chữ, chọn gợi ý, nhập đáp án, dùng sao để mở gợi ý, tính điểm, tính nhiệm vụ sao và ghi nhận tiến độ theo tài khoản. Các chức năng này được triển khai cục bộ, chưa phụ thuộc vào máy chủ hoặc dịch vụ trực tuyến.

\section{Kiểm thử}
\label{section:4.7}

Phần này trình bày khung kiểm thử cho nguyên mẫu. Do báo cáo chưa có dữ liệu kiểm thử chính thức được ghi nhận từ một đợt chạy độc lập, các bảng dưới đây không tự điền trạng thái Passed/Failed. Cột kết quả thực tế được để ở dạng ``Ghi sau khi kiểm thử'' để tránh nhầm lẫn giữa kế hoạch kiểm thử và kết quả đã đo.

\subsection{Kiểm thử chức năng}

\begin{longtable}{|>{\RaggedRight\arraybackslash}p{1.5cm}|>{\RaggedRight\arraybackslash}p{3.0cm}|>{\RaggedRight\arraybackslash}p{4.5cm}|>{\RaggedRight\arraybackslash}p{4.0cm}|}
\caption{Khung kiểm thử chức năng}
\label{table:test_cases}\\
\hline
\textbf{Mã} & \textbf{Chức năng} & \textbf{Dữ liệu/thao tác kiểm thử} & \textbf{Kết quả thực tế} \\ \hline
\endfirsthead
\hline
\textbf{Mã} & \textbf{Chức năng} & \textbf{Dữ liệu/thao tác kiểm thử} & \textbf{Kết quả thực tế} \\ \hline
\endhead
TC01 & Tạo tài khoản & Nhập tên tài khoản mới và xác nhận tạo. Kiểm tra tài khoản xuất hiện trong danh sách và được chọn làm hồ sơ hiện tại. & Ghi sau khi kiểm thử. \\ \hline
TC02 & Chọn tài khoản & Chọn một tài khoản đã tồn tại và xem tiến độ tương ứng ở danh sách chủ đề/màn. & Ghi sau khi kiểm thử. \\ \hline
TC03 & Chọn chủ đề & Chọn một chủ đề trong menu và mở danh sách màn của chủ đề đó. & Ghi sau khi kiểm thử. \\ \hline
TC04 & Chọn màn chơi & Chọn một màn trong chủ đề, kiểm tra hệ thống chuyển sang màn chơi và hiển thị bảng/gợi ý. & Ghi sau khi kiểm thử. \\ \hline
TC05 & Sinh bảng ô chữ & Dùng danh sách từ của một chủ đề để sinh bảng; kiểm tra số từ được đặt và tình trạng bảng rời rạc. & Ghi sau khi kiểm thử. \\ \hline
TC06 & Nhập đáp án đúng & Chọn gợi ý, nhập đáp án đúng có dấu tiếng Việt hoặc có khoảng trắng tương ứng. & Ghi sau khi kiểm thử. \\ \hline
TC07 & Nhập đáp án sai/rỗng & Chọn gợi ý rồi nhập sai hoặc để trống ô nhập. & Ghi sau khi kiểm thử. \\ \hline
TC08 & Dùng gợi ý mở chữ & Chọn một từ chưa hoàn thành, dùng sao để mở một chữ. & Ghi sau khi kiểm thử. \\ \hline
TC09 & Dùng gợi ý mở cả từ & Chọn một từ chưa hoàn thành, dùng sao để mở toàn bộ từ. & Ghi sau khi kiểm thử. \\ \hline
TC10 & Không đủ sao & Đưa tài khoản về trạng thái không đủ sao rồi thử dùng gợi ý. & Ghi sau khi kiểm thử. \\ \hline
TC11 & Tính điểm và nhiệm vụ & Hoàn thành một số từ, kiểm tra điểm, chuỗi đúng liên tiếp, nhiệm vụ và số sao nhiệm vụ. & Ghi sau khi kiểm thử. \\ \hline
TC12 & Lưu tiến độ & Hoàn thành màn, quay về menu và kiểm tra trạng thái sao/hoàn thành của màn. & Ghi sau khi kiểm thử. \\ \hline
TC13 & Tải lại tiến độ & Đóng và mở lại game sau khi đã hoàn thành màn. & Ghi sau khi kiểm thử. \\ \hline
TC14 & Dữ liệu lưu không hợp lệ & Thử với tệp lưu thiếu trường hoặc cấu trúc cũ để kiểm tra cơ chế chuẩn hóa. & Ghi sau khi kiểm thử. \\ \hline
\end{longtable}

\subsection{Kiểm thử giao diện}

\begin{table}[H]
\centering
\small
\begin{tabularx}{\textwidth}{|>{\RaggedRight\arraybackslash}p{3.0cm}|>{\RaggedRight\arraybackslash}X|>{\RaggedRight\arraybackslash}p{3.2cm}|}
\hline
\textbf{Nhóm kiểm tra} & \textbf{Tiêu chí quan sát} & \textbf{Kết quả thực tế} \\ \hline
Menu và tài khoản & Người chơi nhận biết được tài khoản hiện tại, biết cách thêm/chọn/xóa tài khoản. & Ghi sau khi kiểm thử. \\ \hline
Chủ đề và màn chơi & Danh sách chủ đề và màn đủ rõ, không làm người chơi nhầm giữa chọn chủ đề và chọn màn. & Ghi sau khi kiểm thử. \\ \hline
Bảng ô chữ & Ô chữ, gợi ý, vùng nhập và nút thao tác không chồng nhau ở kích thước cửa sổ phổ biến. & Ghi sau khi kiểm thử. \\ \hline
Phản hồi thao tác & Thông báo đúng, sai, thiếu sao, hoàn thành màn hiển thị đủ rõ và không che thao tác chính. & Ghi sau khi kiểm thử. \\ \hline
\end{tabularx}
\caption{Tiêu chí kiểm thử giao diện}
\label{table:ui_test_criteria}
\end{table}

\subsection{Thử nghiệm chơi nhóm nhỏ}

Thử nghiệm chơi nhóm nhỏ, nếu thực hiện, chỉ nên dùng để phát hiện lỗi chức năng và vấn đề giao diện dễ thấy. Do nhóm thử nhỏ thường không đại diện cho toàn bộ người dùng, kết quả không được dùng để khẳng định hiệu quả học tập hay khả năng cải thiện ghi nhớ từ vựng.

\begin{table}[H]
\centering
\small
\begin{tabularx}{\textwidth}{|>{\RaggedRight\arraybackslash}p{4.2cm}|>{\RaggedRight\arraybackslash}p{4.2cm}|>{\RaggedRight\arraybackslash}X|}
\hline
\textbf{Nội dung ghi nhận} & \textbf{Cách ghi nhận} & \textbf{Giá trị hiện tại} \\ \hline
Số người chơi thử & Ghi số người tham gia trong đợt thử nghiệm. & Ghi sau khi kiểm thử. \\ \hline
Số màn thử & Ghi số màn được dùng trong phiên thử. & Ghi sau khi kiểm thử. \\ \hline
Số màn hoàn thành & Ghi số màn người chơi hoàn thành. & Ghi sau khi kiểm thử. \\ \hline
Thời gian hoàn thành trung bình & Tính từ thời gian hoàn thành từng lượt chơi. & Ghi sau khi kiểm thử. \\ \hline
Số lần dùng gợi ý trung bình & Ghi số lần mở chữ hoặc mở cả từ. & Ghi sau khi kiểm thử. \\ \hline
Nhận xét giao diện & Tổng hợp lỗi khó thao tác, khó đọc hoặc phản hồi chưa rõ. & Ghi sau khi kiểm thử. \\ \hline
\end{tabularx}
\caption{Khung ghi nhận thử nghiệm chơi nhóm nhỏ}
\label{table:small_playtest_frame}
\end{table}

\subsection{Hạn chế kiểm thử}

Khung kiểm thử hiện tại chưa thay thế cho một đợt kiểm thử đầy đủ trên nhiều thiết bị và nhiều nhóm người dùng. Một số nội dung như độ khó từ vựng, mức cân bằng sao/gợi ý, khả năng đọc giao diện ở nhiều độ phân giải và hành vi nhập tiếng Việt trên môi trường web cần được kiểm tra thêm. Vì vậy, phần kiểm thử trong chương này chỉ đóng vai trò chuẩn bị tiêu chí đánh giá nguyên mẫu, chưa đưa ra kết luận về hiệu quả học tập.

\section{Triển khai thử nghiệm}
\label{section:4.8}

Ứng dụng hiện được triển khai thử nghiệm dưới dạng project Godot chạy cục bộ. Khi chạy bằng Godot, scene chính \code{Main.tscn} khởi tạo menu và màn chơi. Dữ liệu tiến độ được lưu trong vùng \code{user://} nên không cần máy chủ. Nếu export sang Windows, cần cấu hình export preset trong Godot 4.x và kiểm tra lại khả năng tạo/đọc tệp lưu trên máy người dùng.

Với bản Web, cần kiểm tra thêm khả năng nhập tiếng Việt trong trình duyệt, kích thước giao diện và cơ chế lưu dữ liệu. Báo cáo không khẳng định ứng dụng đã phát hành chính thức; phần triển khai hiện mới ở mức nguyên mẫu chạy cục bộ và chuẩn bị khả năng export thử nghiệm.

\end{document}
